\section{The Formal Core}
\label{sec:formal-core}

\subsection{Inertia As A Composite Quantity}

\begin{definition}[Framework Inertia]
The \textbf{framework inertia} of a state $s$ is
\[
  \Inertia(s) = \alpha R(s) + \beta C(s) + \gamma E(s),
\]
where $R(s)$ is immediate reward intensity, $C(s)$ is cognitive switching cost,
$E(s)$ is environmental reinforcement, and $\alpha,\beta,\gamma > 0$ are
weighting constants.
\end{definition}

\begin{remark}
This is not claimed as a laboratory-complete model. It is a managerial model:
the formula is valuable if it helps locate the main source of transition
failure and guides the choice of intervention.
\end{remark}

\subsection{Activation Barriers}

\begin{definition}[Activation Barrier]
For a transition $s_0 \to s_1$, the \textbf{activation barrier} is
\[
  \Transition(s_0 \to s_1) = V_{\max}(s_0,s_1) - V(s_0),
\]
where $V$ is a potential-like representation of state stability.
\end{definition}

\begin{lemma}[Barrier Amplification]
If the target state requires multiple unresolved micro-decisions, then
$C(s_0)$ rises and therefore $\Transition(s_0 \to s_1)$ rises.
\end{lemma}

\begin{proof}
Unresolved micro-decisions force the agent to perform planning and execution at
the same moment. That enlarges the context-loading burden, which is exactly
what $C(s_0)$ records. Since the barrier depends on the cost of leaving the
current state, the barrier increases.
\end{proof}

\subsection{Volitional Energy}

\begin{definition}[Volitional Energy]
\textbf{Volitional energy} $\Volition(t)$ is the available override capacity
for overcoming the activation barrier at time $t$.
\end{definition}

\begin{axiom}[Finite Override]
There exists a session-dependent upper bound $\Volition_{\max}$ such that
$0 \leq \Volition(t) \leq \Volition_{\max}$.
\end{axiom}

\begin{definition}[Effective Leverage]
The \textbf{effective leverage} of free will at time $t$ is
\[
  L(t) = \frac{\Volition(t)}{\Transition(t)}.
\]
\end{definition}

\begin{theorem}[Decision-Window Advantage]
Free will has maximal operational value at times when the activation barrier is
locally small, not at times when the agent is already deep inside a reinforced
state.
\end{theorem}

\begin{proof}
By definition, leverage increases when the denominator decreases and the
numerator remains fixed. Therefore, if the same quantity of volitional energy is
deployed during a lower-barrier interval, it produces more transition effect.
The strategy implication is immediate: use free will to enter the next state,
not to remain forever in combat against the current one.
\end{proof}

\subsection{A Comparative Table Of Control Variables}

\begin{center}
\renewcommand{\arraystretch}{1.25}
\begin{tabular}{@{}p{0.25\textwidth} p{0.22\textwidth} p{0.4\textwidth}@{}}
\toprule
\textbf{Control Variable} & \textbf{Speed} & \textbf{Interpretation} \\
\midrule
Reduce $R(s)$ & Medium & Remove short-term temptations or delay access \\
Reduce $C(s)$ & Fast & Specify the first step and prepare materials \\
Reduce $E(s)$ & Fast to medium & Change location, posture, tool visibility \\
Increase $\Volition(t)$ & Unreliable & Useful but costly as a primary lever \\
Shift timing & Fast & Use start-of-day or boundary moments \\
\bottomrule
\end{tabular}
\end{center}

\begin{insightbox}
The model shifts strategy from ``be stronger'' to ``lower the barrier before
action is needed.'' This is the core managerial advantage of the framework.
\end{insightbox}

\begin{summarybox}
\begin{itemize}
  \item Inertia is composite, not mystical.
  \item Activation barriers grow when action and planning are forced into the
        same instant.
  \item Decision windows matter because leverage is ratio-based, not purely
        effort-based.
\end{itemize}
\end{summarybox}

